% META: {"id":"1NSI-TC-EX-038","chapitre":"1NSI-TYPES-CONSTRUITS","type_objet":"exercice","capacites":["1NSI-TYPES-CONSTRUITS-C4"],"parcours":1,"competences":["analyser"],"duree_min":8,"mode_creation":"ex_nihilo","sources_inspiration":[],"status":"approved","origin":"needs_review","status_history":["needs_review","audited","accepted","approved"],"release_acceptance":"RELEASE_OWNER_BATCH_ACCEPTANCE","acceptance_closure_digest":"sha256:9b3ccf9a81c5520fbb7e03b7b2d3e7bf2057a02834b6d908bf3be5f49f3b553f","acceptance_receipt_digest":"sha256:4fad86f0282e0fa4e0419cca1ad6379ae7af5a280c04a4a90880f90a1c044242"}
% BEGIN-TRACE
% station = {"nom": "Paris", "temp": 25, "vent": 10}
% for cle in station:
%     print(cle, ":", station[cle])
% EXPECTED
% nom : Paris
% temp : 25
% vent : 10
% END-TRACE
\begin{exercice}{1NSI-TC-EX-038}{1}{8}
On considère le programme suivant :
\begin{python}
station = {"nom": "Paris", "temp": 25, "vent": 10}
for cle in station:
    print(cle, ":", station[cle])
\end{python}
\begin{enumerate}
  \item Donner, sans exécuter, les affichages produits.
  \item Que parcourt la boucle \lstinline{for cle in station} : les clés, les valeurs,
        ou les deux ?
  \item Réécrire la boucle en utilisant \lstinline{station.items()}.
\end{enumerate}
\end{exercice}
